\section{怎样证明描述的确是同一次运行}

结果槽中出现 21 很有吸引力，但它不是充分证据。RAM 可能保留了上一次的 21，操作者也
可能提前把该值写入。要判断这一次运行，必须把准备、交接、执行和返回四个阶段的证据连成
一条不矛盾的链。

\subsection{为每个阶段选择最小而不同的证据}

\begin{longtable}{|L{0.18\linewidth}|L{0.28\linewidth}|L{0.42\linewidth}|}
\hline
\rowcolor{BookTableHead}
要证明的事件 & 使用的证据 & 单独使用时的不足 \\ \hline
本批次已准备完整 & \code{prepared=1}、编号 17、两个校验和 & 只能说明启动前字节通过检查 \\ \hline
任务入口已经提交 & \code{phase=ENTERED}、记录的入口 PC & 不能说明任务走到循环末尾 \\ \hline
结果存储已经发生 & 哨兵被替换为 \code{0x00000015} & 不能单独排除旧结果或其他写入者 \\ \hline
任务已经合作返回 & \code{phase=RETURNED}、编号 17 & 不能替代对结果数值的检查 \\ \hline
没有处理器故障 & \code{halted=0} & 不能证明程序语义正确 \\ \hline
\end{longtable}

启动程序在进入任务前把 \code{phase} 写为 \code{ENTERED}，返回入口把它改为
\code{RETURNED}。两次写入都同时复制准备编号。因为当前只有一个处理器和一个运行程序，
不存在另一个合法写入者；即便如此，仍要把状态字和结果槽共同检查。

\subsection{提交次序给证据建立先后关系}

我们可以从已定义的提交边界推出：

\begin{enumerate}
\item \code{phase=ENTERED} 的写响应成功后，启动程序才提交任务入口 PC；
\item 任务四轮循环结束后，才可能到达结果 \code{STORE}；
\item 结果写响应成功后，任务才取到 \code{RET}；
\item \code{RET} 成功提交后，返回入口才可能写 \code{phase=RETURNED}。
\end{enumerate}

于是同一编号下的 \code{RETURNED} 蕴含控制流已经越过结果存储的成功响应。但它仍不蕴含
结果一定是 21，因为错误的算法也可以合作返回。最终数值要由指令 trace 和输入字节独立
复算。

\begin{center}
\begin{tikzpicture}[node distance=9mm and 11mm]
\node[stage,minimum width=2.5cm] (prepared) {\code{PREPARED}\\记录可核对};
\node[stage,minimum width=2.5cm,right=of prepared] (entered) {\code{ENTERED}\\任务 PC 已提交};
\node[stage,minimum width=2.5cm,right=of entered] (stored) {结果写响应\\成功};
\node[stage,minimum width=2.5cm,right=of stored] (returned) {\code{RETURNED}\\返回入口已执行};
\draw[flow] (prepared) -- (entered);
\draw[flow] (entered) -- (stored);
\draw[flow] (stored) -- (returned);
\end{tikzpicture}
\end{center}
\diagramnote{箭头来自具体指令和提交规则，不是状态名称本身带来的保证。若将来出现并发
写入者或缓存，这条证据链需要重新说明。}

\section{从循环不变量复算结果}

逐条 trace 可以展示每条动态指令；另一个互补方法是找出每轮都保持的关系。设输入数组为
\(a_0,\ldots,a_3\)，刚要开始第 \(k\) 轮时：
\[
\begin{aligned}
\texttt{r0} &= \texttt{data\_base}+4k,\\
\texttt{r1} &= 4-k,\\
\texttt{r4} &= \sum_{i=0}^{k-1}a_i.
\end{aligned}
\]
当 \(k=0\) 时，入口初始化给出 \code{r0=data\_base}、\code{r1=4} 和
\code{r4=0}，关系成立。

假设第 \(k\) 轮开始时关系成立。装入得到 \(a_k\)，加法令累加器成为
\(\sum_{i=0}^{k}a_i\)，指针增加四，计数减少一。因此下一轮开始时：
\[
\texttt{r0}=\texttt{data\_base}+4(k+1),\qquad
\texttt{r1}=4-(k+1),
\]
不变量继续成立。当第四轮结束时 \code{r1=0}，分支退出，累加器为
\[
7+(-2)+11+5=21.
\]

这个推导依赖三项此前已具体建立的事实：\code{r1} 初值确为四、每个元素恰为四字节、
回边指向装入而不指向清零。若任何一项变化，证明必须随之修改。

\subsection{32 位加法是否会溢出}

本例的所有部分和都位于有符号 32 位范围，故位型加法与数学整数加法得到相同解释。一般
输入可能溢出。本章处理器的 \code{ADD} 保留低 32 位，不自动停止：
\[
\texttt{0x7fffffff}+1=\texttt{0x80000000}.
\]
若把结果解释为有符号数，它是 \(-2147483648\)。因此“求和”在当前指令契约中准确含义是
32 位模加法；本例恰好无需区分。需要报告算术溢出的程序必须检查标志或使用更宽表示，不能
从自然语言“相加”推断硬件会代劳。

\section{用反事实检查准备协议的每一道边界}

好的协议不仅要说明成功路径，还应能回答“少做一步会怎样”。下面每个实验只改变一个准备
动作。

\subsection{没有先清零旧的提交标志}

RAM 中仍有编号 16 的完整记录，操作者直接开始覆盖代码，尚未写完就释放复位。若启动程序
只检查 \code{prepared=1}，它可能扫描一半新、一半旧的代码。加入准备编号和校验和后，
最迟在完整性检查处停止；而正确操作在第一步清零标志，可以更早明确“尚未完成”。

\subsection{只检查代码首地址，不检查长度}

令 \code{code\_base=0x001ffff0}、\code{code\_bytes=32}。首地址位于 RAM，但区域末端
越过 \code{0x001fffff}。若启动程序逐字节复算校验和，第十七个字节已经没有目标，互连
返回 \code{UNMAPPED}。预先做不回绕的范围检查能在任何扫描前拒绝该记录，并把原因定位到
字段关系，而不是一次偶然的访存失败。

\subsection{代码与输入区域重叠}

若 \code{data\_base=0x00100010}，后两个输入整数会覆盖计数更新和分支指令。两个区域各自
都位于 RAM，两个校验和也可能分别按被覆盖后的内容计算成功；只有跨区域重叠检查能指出
同一批字节被赋予了互不相容的角色。

\subsection{先改变 PC，再写完任务寄存器}

假设交接代码先提交 \code{PC=0x00100000}，下一周期才准备 \code{r0}。第一条任务指令
清零 \code{r4} 后，第二条立即读取旧 \code{r0}，可能访问任意地址。入口 PC 与参数必须
属于同一个交接提交边界，或者硬件必须保证提交后才开始取指。本章采用前者。

\subsection{结果写入成功，但返回栈字损坏}

四轮计算和结果 \code{STORE} 均可成功，随后 \code{RET} 读到未对齐地址并使处理器停止。
此时结果槽为 21，\code{phase} 仍为 \code{ENTERED}，故障原因为返回目标非法。这个例子
说明“结果已产生”与“任务已合作返回”是两个不同事件；状态设计需要同时容纳二者。

\begin{longtable}{|L{0.22\linewidth}|L{0.25\linewidth}|L{0.41\linewidth}|}
\hline
\rowcolor{BookTableHead}
单点变化 & 预期停止位置 & 由哪项状态区分 \\ \hline
提交标志未清零 & 编号或校验检查 & \code{phase} 与 \code{detail} \\ \hline
代码区域越过 RAM & 范围检查 & 失败字段偏移 \\ \hline
代码与输入重叠 & 区域关系检查 & 两个区域编号 \\ \hline
入口寄存器非原子建立 & 任务早期访存 & \code{fault\_pc} 与数据地址 \\ \hline
返回栈字损坏 & \code{RET} & 结果已变、状态仍为 \code{ENTERED} \\ \hline
\end{longtable}

\section{重复运行为什么不能只按一次复位}

第一次任务返回后，RAM 至少有三处变化：结果槽已从哨兵变为 21，准备状态变为
\code{RETURNED}，任务也可能改写自己的栈。若操作者只复位处理器，启动程序看到的仍是
编号 17 和旧结果。

本章规定每次运行采用新的准备编号，并重新执行以下动作：

\begin{enumerate}
\item 保持复位，先把 \code{prepared} 清零；
\item 写入新的编号，恢复输入、结果哨兵和初始栈字；
\item 必要时重写代码并读回全部声明区域；
\item 写入新的校验和，最后重新提交；
\item 释放复位。
\end{enumerate}

即使代码没有变化，也不能省略恢复数据与结果槽，因为程序语义可能依赖初值。若希望快速
重复运行而不重写不变代码，需要进一步区分只读模板与每次运行的可变状态，并由机器自己
完成复制；这已经是后续演进，不应暗含在本章的人工步骤中。

\subsection{复位、重新准备和重新运行是三个不同动作}

\begin{longtable}{|L{0.20\linewidth}|L{0.28\linewidth}|L{0.40\linewidth}|}
\hline
\rowcolor{BookTableHead}
动作 & 确定改变的状态 & 不保证改变的状态 \\ \hline
复位 & PC、控制状态、\code{halted} & RAM 中的代码、输入、结果与记录 \\ \hline
重新准备 & 声明的 RAM 区域与准备编号 & 处理器当前 PC 和内部执行状态 \\ \hline
重新运行 & 从固定入口推进并产生新效果 & 断电后的保存、其他机器上的结果 \\ \hline
\end{longtable}

所以正确顺序是先保持复位，再准备，最后释放复位。运行中直接让夹具改代码，会产生一份
本章没有定义的交错历史。

\section{当前机器的责任边界}

经过完整走读，我们可以把每个部件承担的责任写得足够窄。

\begin{longtable}{|L{0.18\linewidth}|L{0.33\linewidth}|L{0.37\linewidth}|}
\hline
\rowcolor{BookTableHead}
部件或参与者 & 它负责保证 & 它不负责保证 \\ \hline
操作者与夹具 & 复位期间按顺序写入并读回声明字节 & 运行中自动提供新程序 \\ \hline
启动存储 & 复位后提供不变的检查与返回指令 & 保存任务运行产生的结果 \\ \hline
RAM & 对成功事务保存和返回相应字节 & 理解代码、输入、栈等含义 \\ \hline
地址互连 & 每个地址选择唯一目标并返回错误 & 阻止任务访问别的 RAM 区域 \\ \hline
处理器 & 按指令提交规则更新 PC、寄存器与访存效果 & 保证任务返回或算法正确 \\ \hline
启动程序 & 验证人工记录并建立一次入口状态 & 并发运行、自动装入或故障恢复 \\ \hline
任务 & 按约定读取参数、写结果并合作返回 & 保护启动程序工作区 \\ \hline
\end{longtable}

这张表没有额外创造一个新层次。它只是防止把某个部件的局部承诺扩大成整机承诺。例如 RAM
能够正确保存写入字节，不等于任务不能写错地址；启动程序检查入口合法，也不等于任务运行
中只能在代码区取指。

\subsection{首次运行已经得到的三个重要分界}

第一，\emph{字节存在}与\emph{可以开始执行}不同。代码、数据、栈和入口状态共同通过检查
以后，PC 的提交才建立开始事件。

第二，\emph{候选效果}与\emph{已提交效果}不同。装入等待响应时不能提前写目标寄存器，
存储失败时不能留下半个新整数，交接失败时不能让任务看到半套参数。

第三，\emph{结果出现}与\emph{控制权返回}不同。结果写入发生在 \code{STORE} 提交，
控制权返回发生在 \code{RET} 提交；状态记录让两者可以分别观察。

这些分界会在后面的机器演进中反复出现，但后续章节必须从新的具体问题重新推导它们的
用途，而不能只复述三个口号。

\subsection{当前机器仍然不能做的事情}

\begin{itemize}
\item 操作者必须逐字节或逐块写 RAM，程序越大，准备时间越长；
\item 每次换程序都要人工计算地址、校验和、栈字和记录字段；
\item 运行中的任务独占处理器，除故障停止和整机复位外，没有强制收回途径；
\item 所有 RAM 地址对任务可见，初始分区没有运行期保护；
\item RAM 断电后不提供本章可依赖的长期保存；
\item 机器一次只有一组任务寄存器，没有第二个可暂停后继续的现场。
\end{itemize}

下一步最直接的困难不是同时运行很多任务，而是人工装入本身。假设程序由几万字节组成，
继续要求操作者借助夹具写入每个地址，既慢又容易把地址、长度或字节顺序抄错。要让机器从
一条外部字节通道自动接收同样的代码与数据，就必须新增接收硬件，并让启动程序承担接收、
校验和最后交接。这个问题将单独构成下一章。
